\Article{10}
Als Beispiel wollen wir die Formel auf die magnetischen Elemente für
\begin{center}
\begin{tabular}{l|l|l|l}
Göttingen&$\delta^0=18^0\ 38'$&$i^0=67^0\ 56'$&$\psi^0=1{,}357$\\
Mailand&$\delta'=18\ 33$&$i'=63\ 49$&$\psi'=1{,}294$\\
Paris&$\delta''=22\ 4$&$i''=67\ 24$&$\psi''=1{,}348$
\end{tabular}
\end{center}
anwenden. Damit erhält man
\[
\begin{aligned}
\omega^0&=0{,}50980\\
\omega'&=0{,}57094\\
\omega''&=0{,}51804
\end{aligned}
\]
Legt man die geografische Lage
\begin{center}
\begin{tabular}{l|l|l}
Göttingen&$51^0\ 32'$ Breite&$9^0\ 58'$ Länge von Greenwich\\
Mailand&$45\ 28$&$9\ 9$\\
Paris&$48\ 52$&$2\ 21$
\end{tabular}
\end{center}
zugrunde und rechnet lediglich wie auf der Kugelfläche, so erhält man
\[
\begin{aligned}
\left.\begin{aligned}(01)&=5^0\ 11'\ 31''\\(10)&=184\ 35\ 35\end{aligned}\right\}&\quad P^0P'=6^0\ 5'\ 20''\\[4pt]
\left.\begin{aligned}(12)&=128\ 47\ 31\\(21)&=303\ 48\ 1\end{aligned}\right\}&\quad P'P''=5\ 44\ 6\\[4pt]
\left.\begin{aligned}(20)&=238\ 20\ 20\\(02)&=64\ 10\ 12\end{aligned}\right\}&\quad P^0P''=5\ 32\ 4
\end{aligned}
\]
\par\SourceLeaf{137}
Setzt man diese Werte und die obigen für $\delta^0,\delta',\delta''$ in unsere Gleichung ein und drückt die Abstände in Bogensekunden aus, so erhält man
\[
0=17556\omega^0+2774\omega'-20377\omega''
\]
oder
\[
\omega''=0{,}86158\omega^0+0{,}13613\omega'
\]
Aus den beobachteten horizontalen Intensitäten in Göttingen und Mailand folgt damit für Paris $\omega''=0{,}51696$, in fast genauer Übereinstimmung mit dem beobachteten Wert $0{,}51804$.

Übrigens sieht man leicht, dass sich die obige Formel unmittelbar durch die geografischen Längen und Breiten der Orte ausdrücken lässt, wenn man die Abstände $P^0P'$ usw. durch ihre Sinus ersetzen will.

